\chapter*{VI/16 Legacy-Problem Audit}
\addcontentsline{toc}{chapter}{VI/16 Legacy-Problem Audit}

\paragraph{Audit rule.}
Every substantive legacy problem assigned to Volume VI must have exactly one canonical
destination.  For VI/16 the audit started from the volume-wide ledger, then checked all three
exactness source streams and their repeated theorem/solution variants.

\paragraph{Principal preserved corpus problems.}
\begin{enumerate}
\item \textbf{AG1 Exercise 6} --- the full theorem that a sequence of sheaves is exact if and
only if every stalk sequence is exact.  Action: \textbf{INCLUDE/MERGE}; canonical destination
VI.16.P1.
\item \textbf{AG2 Problem 7} --- a morphism of sheaves is an isomorphism if and only if it is
injective and surjective, with an explicit germ-lifting, overlap-agreement, and gluing proof.
Action: \textbf{INCLUDE/MERGE}; canonical destination VI.16.P2.
\end{enumerate}

\paragraph{Merged parallel source material.}
The preferred \texttt{sheaves\_notes\_reorganized.tex} exactness/stalkwise section contains the
same central theorem, together with the exponential sequence, de Rham examples, and the warning
that global sections need not preserve surjectivity.  These are merged into the theory, examples,
and VI.16.P1 rather than duplicated as additional legacy problems.  AG1 also repeats the
stalkwise criterion later as summary/theorem material; AG2 repeats the isomorphism criterion as a
theorem after Problem 7.  Those repetitions are variants, not independent mathematical problems.

\paragraph{Reserved neighboring corpus problem.}
AG2 Problem 9 concerns left exactness of $\Gamma(U,-)$ and the later cohomological consequences.
VI/16 explains the underlying phenomenon but does \emph{not} consume that numbered problem.  The
ledger keeps it reserved for VI/48, where the global-section/cohomology sequence is developed at
full depth.

\paragraph{Coverage conclusion.}
The three migration scopes assigned to VI/16 --- the preferred notes exactness descendants, AG1
exactness descendants, and AG2 exactness descendants --- have all been instance-audited.  No
additional unique numbered exactness problem was found in these scopes.  Therefore the chapter's
two legacy dossiers are the complete numbered VI/16 corpus recovery; the remaining dossiers are
new synthesis problems and are labeled honestly as such.
